\section{A Green estimate in the finite map law}
\label{sec:green}

Uniform approximation of energies concerns continuous functions.  To
control the walk clock, we also need to show that bounded sources of
small stationary mass have small potentials.  We prove this directly
under the finite map law.  The argument combines a spanning-forest
formula for the inverse Laplacian with the two tree contours in
Mullin's encoding.

Write $N=2n$.  Let $H_n$ be the positive graph Laplacian on $V_n$,
using counting measure and unit conductances, so that
\[
 \langle u,H_nu\rangle=\en_n(u)
       =\sum_{ab\in E(M_n)}(u(a)-u(b))^2.
\]
There is one term for each unoriented edge, and loops contribute zero.
For a real function $f$ with $\mu_n(f)=0$, define $G_nf$ by
\begin{equation}\label{gr:green-definition}
 H_nG_nf=\mu_n f,
 \qquad \mu_n(G_nf)=0.
\end{equation}
Here $\mu_nf$ on the right denotes the vector with entry
$\mu_n(v)f(v)$ at $v$.  Connectedness gives existence and uniqueness.

\begin{proposition}[Small-mass Green estimate]
\label{gr:green-bound}
For each fixed $\rho\in(0,1/2)$ there are nonnegative random variables
$A_n$, measurable functions of $M_n$, such that
\begin{equation}\label{gr:green-main}
 \sup_{n\geq1}\E A_n<\infty,
 \qquad
 \|G_nf\|_\infty\leq A_n\,\mu_n(|f|)^\rho
\end{equation}
simultaneously for all real functions $f:V_n\to\R$ satisfying
$\|f\|_\infty\leq1$ and $\mu_n(f)=0$.
The expectation is under the map marginal of the uniform rooted
map--spanning-tree pair law.  In particular, $(A_n)$ is tight.
\end{proposition}

We first identify the exact inverse formula.  We then estimate its
terms by grouping tree cuts according to the amount of tour time on
their smaller side.

\subsection{The inverse Laplacian as a sum over cuts}

The next identity is a consequence of the matrix-tree theorem and
its forest expansions; see the all-minors theorem of \cite{Chaiken}
and \cite[Chapter~4]{LP} for electrical networks and spanning trees.
We give the determinant argument to specify the multiplicities and
centering used here.

\begin{lemma}[Forest representation]
\label{gr:forest}
Let $\mathcal G$ be a finite connected multigraph, with positive
probability measure $\mu$ on its vertex set $V$, and let $H$ be its
unit-conductance graph Laplacian.  Suppose $q:V\to\R$ satisfies
$\sum_{v\in V}q(v)=0$, and let $u$ solve
$Hu=q$, $\mu(u)=0$.  For a uniform spanning tree $T$ of $\mathcal G$
and an edge $e\in T$, choose either component $S_e$ of $T\setminus e$,
and let $b_e$ be the number of graph edges between $S_e$ and
$V\setminus S_e$, with multiplicity.  Then
\begin{equation}\label{gr:forest-identity}
 u(x)=\E_T\sum_{e\in T}
       \frac{q(S_e)\bigl(\1_{\{x\in S_e\}}-\mu(S_e)\bigr)}{b_e},
 \qquad x\in V,
\end{equation}
where $q(S)=\sum_{v\in S}q(v)$.  The value does not depend on the
choice of component for any $e$.
\end{lemma}

\begin{proof}
The one-vertex case is immediate, so suppose $m:=|V|\geq2$.
Orient the non-loop edges arbitrarily and delete the row of a fixed
vertex $o$ from the incidence matrix.  Denote the resulting matrix by
$B$, and write $\widehat q$ for $q$ with its $o$-coordinate deleted.
The reduced Laplacian is $\widehat H=BB^{\mathsf T}$.  If $\tau$ is
the number of spanning trees, the matrix-tree theorem and the
determinant lemma give
\begin{equation}\label{gr:rank-one}
 \det(\widehat H+z\widehat q\widehat q^{\mathsf T})
       =\tau\bigl(1+z\widehat q^{\mathsf T}
                            \widehat H^{-1}\widehat q\bigr).
\end{equation}
For $z\geq0$, apply Cauchy--Binet to the matrix obtained by appending
$\sqrt z\widehat q$ to $B$.  A term contributing to the coefficient
of $z$ selects $m-2$ incidence columns and the source column.
The incidence columns are independent exactly when their edges form
a spanning forest with two components.  Indeed, a cycle gives a
linear dependence, whereas an acyclic graph with $m-2$ edges has
exactly two components, counting isolated vertices.

For such a forest $F$, let $S_F$ be its component not containing $o$.
Its corresponding determinant has absolute value $|q(S_F)|$.
To see this, successively eliminate leaf rows and their incident
forest columns.  Eliminating a leaf adds its source entry to its
neighbor's entry.  The component containing $o$ disappears into the
deleted row, and the last entry in the other component is
$\sum_{v\in S_F}q(v)$.  Consequently, comparison of the coefficients
of $z$ in \eqref{gr:rank-one} yields
\begin{equation}\label{gr:two-forest-quadratic}
 \widehat q^{\mathsf T}\widehat H^{-1}\widehat q
       =\frac1\tau\sum_F q(S_F)^2.
\end{equation}

Let $H^+$ be the inverse of $H$ on counting-mean-zero vectors,
extended by zero on constants.  Solving the equation with value zero
at $o$ shows that the left side of
\eqref{gr:two-forest-quadratic} equals $q^{\mathsf T}H^+q$.
Polarization therefore gives, for any two vectors $q,r$ of total
sum zero,
\begin{equation}\label{gr:two-forest-bilinear}
 q^{\mathsf T}H^+r=\frac1\tau\sum_F q(S_F)r(S_F).
\end{equation}
Replacing $S_F$ by its complement changes the sign of each factor,
so the product is independent of this choice.

Every two-component forest $F$ has exactly $b(F)$ extensions to a
spanning tree: add any one of the graph edges across its cut.
Thus, when all pairs $(T,e)$ are counted, $F=T\setminus e$ occurs
exactly $b(F)$ times.  This explains the denominator in
\eqref{gr:forest-identity}.  Apply \eqref{gr:two-forest-bilinear}
with $r=\delta_x-\mu$.  Its left side is
\[
 (H^+q)(x)-\mu(H^+q)=u(x),
\]
and its right side, regrouped by $(T,e)$, is the asserted expression.
Loops enter neither the reduced incidence matrix nor a cut; parallel
edges are counted separately throughout.
\end{proof}

\subsection{What a tree cut records in the tour}

We use the finite form of Mullin's correspondence \cite{Mullin}, with
the incidence conventions of \cite[Section~``Independent-walk encoding
and local quantum surfaces'']{MetricMaps}.  The $N$ steps form a
quadrant excursion whose horizontal and vertical coordinates are
denoted by $L$ and $R$.  Horizontal steps traverse the selected
primal tree $T$; vertical steps traverse the complementary dual tree
$T^*$.  In each coordinate an up-step is matched with the down-step
that first returns to its initial height.  A matched horizontal pair
represents a tree edge.  A matched vertical pair represents a
non-tree map edge between the primal vertices at its two step starts.
Let $v_i$ be the primal vertex at step start $i$, with all indices
read cyclically modulo $N$.

\begin{lemma}[Tour mass and cut size]
\label{gr:tour-cut}
For each $v\in V_n$,
\begin{equation}\label{gr:tour-degree}
 \#\{0\leq i<N:v_i=v\}=\deg_{M_n}(v).
\end{equation}
For $e\in T$, the step starts visiting either component of
$T\setminus e$ form a cyclic interval.  Choose a component $S_e$
whose interval $I_e$ has length $\ell_e\leq N/2$.  If the lengths
agree, choose the component not containing $v_0$.  Then
\begin{equation}\label{gr:tour-cut-identities}
 \mu_n(S_e)=\frac{\ell_e}{N},
 \qquad b_e=1+D_e^*,
\end{equation}
where $D_e^*$ is the distance in $T^*$ between the dual vertices at
the two horizontal traversal steps of $e$.
\end{lemma}

\begin{proof}
Associate each step start with the departing step.  The two
horizontal traversals of a tree edge depart once from each endpoint,
and hence count its two incidences.  The two steps of a vertical
matched pair count the two incidences of its non-tree edge.  When
that edge is a loop, both are assigned to the same vertex.
Every step belongs to one such pair, proving
\eqref{gr:tour-degree}; this is also the tour-measure lemma of
\cite[Section~``From tree excursions to the measured quantum
sphere'']{MetricMaps}.

The primal vertex changes only along horizontal steps.  It passes
between the two components of $T\setminus e$ exactly at the two
traversals of $e$.  Thus each side occupies one cyclic interval, and
\eqref{gr:tour-degree} proves the mass identity.

There is exactly one tree edge across this cut, namely $e$.
A non-tree edge crosses it precisely when exactly one step of its
vertical matched pair belongs to $I_e$.  Ignore horizontal holds
and follow the resulting dual-tree contour segment.  An edge of
$T^*$ is traversed once on this segment exactly when its matched
steps lie on opposite sides of $I_e$; every other dual edge is
traversed zero or twice.  In a tree the edges traversed an odd number
of times by a path are exactly the edges on the simple path between
its endpoints.  These endpoints are the dual vertices at the two
horizontal traversals of $e$, since horizontal steps do not change
the dual vertex.  There are therefore $D_e^*$ non-tree edges across
the cut, giving $b_e=1+D_e^*$.
\end{proof}

Hereafter expectations involving both $T$ and the map use the uniform
pair law.  The sequence $(v_i)$ is determined by the horizontal
contour and the locations of the horizontal slots.  Thus the cut
intervals, their lengths, and the smaller-side convention in
Lemma~\ref{gr:tour-cut} are determined without inspecting the shape
of the vertical contour.  By \eqref{gr:tour-degree}, the same is true
of the vertex masses.

The conditional structure of the finite law is useful.  Let $J$
be the number of horizontal steps.  Conditional on $J$, the horizontal
slots form a uniform $J$-subset of $\{0,\ldots,N-1\}$.  Conditional
on these slots, the two coordinate strings, with holds deleted, are
independent uniform Dyck excursions of lengths $J$ and $N-J$.
In fact, writing $C_k=(k+1)^{-1}\binom{2k}{k}$ for the Catalan
numbers, the number of quadrant excursions with $J=2k$ is
\begin{equation}\label{gr:schedule-count}
 \binom{2n}{2k}C_kC_{n-k},
 \qquad
 \sum_{k=0}^n\binom{2n}{2k}C_kC_{n-k}=C_nC_{n+1}.
\end{equation}
The last equality is Vandermonde's identity after writing out the
Catalan numbers.  These counts also prove the stated conditional
independence.  Catalan estimates and a binomial tail bound give
absolute $C,c>0$ such that
\begin{equation}\label{gr:count-tail}
 \P\{J\notin[N/4,3N/4]\}\leq C(N+1)^3e^{-cN}.
\end{equation}
The same schedule calculation is recorded in the excursion discussion
of \cite[Section~``From tree excursions to the measured quantum
sphere'']{MetricMaps}.  It is an exact calculation in the fixed-size
law.

\subsection{Distances in a uniform contour tree}

The cut denominator in \eqref{gr:tour-cut-identities} is a dual-tree
distance.  The next estimates quantify its typical size and also
control how many primal cuts can surround a given tour time.

\begin{lemma}[Contour distance estimates]
\label{gr:dyck}
Let $\mathcal T_s$ be a uniform rooted plane tree with $s/2$ edges,
where $s\geq2$ is even, and write $w_i$, $i\in\Z/s\Z$, for its
cyclic contour vertices.  If $i,j$ are deterministic with
$0\leq j\leq s$, set
\[
 D=\dist_{\mathcal T_s}(w_i,w_{i+j}),
 \qquad d=1+\min(j,s-j).
\]
There are absolute constants $C,c>0$ such that, for every integer
$a\geq0$,
\begin{equation}\label{gr:dyck-density}
 \P(D=a)\leq C(a+1)^2d^{-3/2}e^{-ca^2/d}.
\end{equation}
Consequently, for every fixed $p>0$,
\begin{equation}\label{gr:dyck-moments}
 \|(1+D)^{-1}\|_{L^2}\leq C d^{-1/2},
 \qquad \|D\|_{L^p}\leq C_p\sqrt d.
\end{equation}
For every deterministic cyclic segment $K$ from corner $i$ through
corner $i+L$, where $1\leq L\leq s$,
\begin{equation}\label{gr:contour-diameter}
 \bigl\|\diam_{\mathcal T_s}\{w_{i+r}:0\leq r\leq L\}\bigr\|_{L^4}
       \leq C\sqrt L.
\end{equation}
\end{lemma}

\begin{proof}
Rerooting a plane tree at its $i$th contour corner is a bijection on
rooted plane trees with $s/2$ edges.  Its inverse reroots by $s-i$
corners.  Thus the rerooted tree is uniform, and $D$ has the law of
the height at time $j$ of a uniform Dyck excursion of length $s$.

Let $p_l(a)$ be the probability that a simple symmetric walk started
at zero stays nonnegative for $l$ steps and finishes at $a$.
Reflection and the binomial estimate give
\begin{align}
 p_l(a)
 &=2^{-l}\left[
       \binom{l}{(l+a)/2}-\binom{l}{(l+a+2)/2}\right],
       \label{gr:killed-kernel}\\
 p_l(a)&\leq C(a+1)(l+1)^{-3/2}
                         e^{-ca^2/(l+1)},\nonumber
\end{align}
with inadmissible binomial coefficients interpreted as zero.
For even $s$, $p_s(0)$ is comparable to $(s+1)^{-3/2}$.
The Markov property and time reversal of the second segment give
\[
 \P(D=a)=\frac{p_j(a)p_{s-j}(a)}{p_s(0)}.
\]
Inserting \eqref{gr:killed-kernel} proves
\eqref{gr:dyck-density}.  Summing it after multiplication by
$(a+1)^{-2}$ or by $a^p$ proves
\eqref{gr:dyck-moments}.  In the first case the remaining Gaussian
sum is at most $C\sqrt d$, so the second moment is at most $C/d$.

For completeness, the diameter estimate follows by dyadic chaining.
Reroot at the deterministic first corner of $K$ and parametrize its
steps by the integers in $[0,L]$.  At refinement level $l$, use the
times $\lfloor kL/2^l\rfloor$, $0\leq k\leq2^l$, and stop at
$l=\lceil\log_2L\rceil$.  These grids are nested and the last one
contains every integer in $[0,L]$.  Include the initial increment
from $0$ to $L$.  For $l\geq1$, there are $O(2^l)$ new-to-parent
increments at level $l$, each spanning at most $C L2^{-l}$ steps.
By the fourth-moment estimate and
$\|\max_k Z_k\|_{L^4}^4\leq\sum_k\|Z_k\|_{L^4}^4$, the maximum
tree distance across these increments has $L^4$ norm at most
\[
 C2^{l/4}\sqrt{L2^{-l}}=C\sqrt L\,2^{-l/4}.
\]
Every contour vertex can be joined to the first vertex through its
ancestors in the dyadic grids.  The triangle inequality and the
summability of $2^{-l/4}$ bound the maximum distance from that first
vertex by $C\sqrt L$ in $L^4$.  The diameter is at most twice this
maximum.  For $L=1$ the bound follows directly from the single edge.
\end{proof}

\subsection{The sum of reciprocal cut sizes at one scale}

For a dyadic integer $M$ with $1\leq M\leq N/2$, let
\[
 \mathcal E_M=\{e\in T:M\leq\ell_e<2M\},
 \qquad
 U_{M,t}=\sum_{\substack{e\in\mathcal E_M\\t\in I_e}}\frac1{b_e},
 \qquad t\in\Z/N\Z.
\]
The intervals in this sum come from the horizontal tree.  Their
denominators are governed by the independent vertical tree once the
schedule and the horizontal tree have been fixed.

\begin{lemma}[One-scale estimate]
\label{gr:one-scale}
There is an absolute constant $C$ such that, for every $n$, every
dyadic $M\leq N/2$, and every deterministic tour time $t$,
\begin{equation}\label{gr:Ut}
 \E U_{M,t}^2\leq C.
\end{equation}
\end{lemma}

\begin{proof}
First condition on a count $J\in[N/4,3N/4]$.  Consider all cyclic
intervals containing $t$ whose lengths lie in $[M,2M)$ and are at
most $N/2$.  There are at most $C(M+1)^2$ such intervals.  For each
one, a uniform schedule puts at least $c_0M$ vertical steps in both
the interval and its complement, except with probability
$Ce^{-cM}$, for suitable absolute $c_0,c>0$.  Indeed, the vertical
density is at least $1/4$, both sets have at least $M$ slots, and
the hypergeometric lower-tail bound applies.  Taking a union bound,
let $\mathcal B_{M,t}$ be the event that this property fails for
some interval.  Uniformly over the balanced counts,
\begin{equation}\label{gr:schedule-tail}
 \P(\mathcal B_{M,t}\mid J)
       \leq C(M+1)^2e^{-cM}.
\end{equation}

Fix a schedule outside $\mathcal B_{M,t}$ and then condition on the
horizontal Dyck excursion.  This fixes all intervals in $U_{M,t}$.
For each one, its dual endpoints are separated in the vertical
contour by the number $j$ of vertical steps in that interval.
Both $j$ and $(N-J)-j$ are at least $c_0M$.  Lemmas
\ref{gr:tour-cut} and \ref{gr:dyck}, followed by the triangle
inequality in $L^2$ of the remaining vertical excursion, imply
\begin{equation}\label{gr:vertical-bound}
 \|U_{M,t}\|_{L^2(R\mid\text{schedule},L)}
       \leq\frac{C}{\sqrt M}\,K_{M,t},
 \qquad
 K_{M,t}:=\#\{e\in\mathcal E_M:t\in I_e\}.
\end{equation}

We next bound this number of horizontal cuts.  Any two components
$S_e,S_{e'}$ containing $v_t$ are nested unless their union is all
of $V_n$.  This elementary tree property follows by considering
the components cut off from $v_t$: they are either nested or
disjoint.  The union alternative cannot occur for the chosen
smaller sides, since
\[
 \mu_n(S_e\cup S_{e'})
 =\mu_n(S_e)+\mu_n(S_{e'})-\mu_n(S_e\cap S_{e'})<1.
\]
Here each of the first two terms is at most $1/2$, and the
intersection contains $v_t$, whose mass is positive.
Thus all components counted by $K_{M,t}$ are nested.

If the family is nonempty, let $S$ be its largest component and
choose a contour vertex immediately outside its interval.  It has
a representative at cyclic distance at most $2M+2$ from $t$.
The tree path from $v_t$ to that vertex crosses all $K_{M,t}$ cut
edges.  Hence $K_{M,t}$ is at most the diameter of the horizontal
contour image of the deterministic window of tour times within
$2M+2$ of $t$.  After the schedule has been fixed and holds deleted,
this is a deterministic contour segment with at most $4M+4$
horizontal steps.  If the window covers the entire tour, the full
horizontal contour instead has length $J\leq N\leq4M+4$.
A segment with no horizontal steps has diameter zero; otherwise
Lemma~\ref{gr:dyck} applies.  Uniformly in the schedule, we obtain
\begin{equation}\label{gr:horizontal-count}
 \E\bigl[K_{M,t}^2\mid\text{schedule}\bigr]\leq CM.
\end{equation}
Combining \eqref{gr:vertical-bound} and
\eqref{gr:horizontal-count} proves the required bound on schedules
outside $\mathcal B_{M,t}$ with balanced counts.

For the exceptional schedules, use a deterministic bound.
Every interval counted by $K_{M,t}$ has its two delimiting tree
traversals within cyclic distance $2M+2$ of $t$.  Distinct tree
edges have distinct traversal occurrences, so
$K_{M,t}\leq C(M+1)$.  Since $b_e\geq1$, the same bound holds
for $U_{M,t}$.  By \eqref{gr:schedule-tail}, the contribution of
these schedules to its second moment is at most
$C(M+1)^4e^{-cM}$, uniformly bounded in $M$.
Finally, $U_{M,t}\leq N$ always.  The unbalanced counts contribute
at most $CN^2(N+1)^3e^{-cN}$ by \eqref{gr:count-tail}, also
uniformly bounded.  This proves \eqref{gr:Ut} under the full law.
\end{proof}

\subsection{Summing the scales}

We now return to the inverse formula.  Lemma~\ref{gr:one-scale}
controls the cut sum at a fixed time; a grid converts this into a
uniform bound over vertices.

\begin{proof}[Proof of Proposition~\ref{gr:green-bound}]
The one-vertex case has $G_nf=0$, so set $A_n=0$ there.  Otherwise,
for each dyadic $M\leq N/2$, choose a deterministic cyclic grid
$\Gamma_M$.  For $M\geq8$, take its gaps at most $M/4$, its
cardinality at most $CN/M$, and its distinct points separated by
at least $cM$.  For $M<8$, use all step indices; the same
cardinality and separation bounds hold with absolute constants.

An interval of length in $[M,2M)$ containing a fixed time $t$
contains a grid point within cyclic distance $2M$ of $t$.  There
are only a bounded number of grid points in this neighborhood.
Consequently,
\begin{equation}\label{gr:grid-comparison}
 \sup_t U_{M,t}\leq C\max_{t\in\Gamma_M}U_{M,t}.
\end{equation}
Every visit to a given vertex belongs to the same side of a tree
cut.  Also, counting the visits in each interval gives
\[
 \sum_{e\in\mathcal E_M}\frac{\mu_n(S_e)}{b_e}
       =\frac1N\sum_{t=0}^{N-1}U_{M,t}.
\]
It follows that
\begin{align}
 W_M&:=\sup_{x\in V_n}\sum_{e\in\mathcal E_M}
        \frac{\1_{\{x\in S_e\}}+\mu_n(S_e)}{b_e}
       \leq C\max_{t\in\Gamma_M}U_{M,t},
          \label{gr:WM}\\
 \E W_M&\leq C\left(\sum_{t\in\Gamma_M}
                                    \E U_{M,t}^2\right)^{1/2}
       \leq C\sqrt{N/M}.\label{gr:WM-mean}
\end{align}
No independence among the grid variables is needed here.

Fix $f$ as in the proposition and write $m_f=\mu_n(|f|)$.
For $e\in\mathcal E_M$, the source mass obeys
\begin{equation}\label{gr:source-interpolation}
 \left|\sum_{v\in S_e}\mu_n(v)f(v)\right|
 \leq\min(2M/N,m_f)
 \leq(2M/N)^{1-\rho}m_f^\rho.
\end{equation}
Define, for the realized pair $(M_n,T)$,
\[
 Y_n(M_n,T)=\sum_{\substack{M\leq N/2\\M\text{ dyadic}}}
                   (2M/N)^{1-\rho}W_M.
\]
By \eqref{gr:WM-mean},
\begin{equation}\label{gr:Y-mean}
 \E Y_n\leq C_\rho
       \sum_{\substack{M\leq N/2\\M\text{ dyadic}}}
                        (M/N)^{1/2-\rho}
       \leq C_\rho.
\end{equation}
The last bound is a geometric sum, using $\rho<1/2$.

Conditional on $M_n$, the selected tree is uniform.  Define
\begin{equation}\label{gr:An}
 A_n(M_n)=\E\bigl[Y_n(M_n,T)\mid M_n\bigr].
\end{equation}
Apply Lemma~\ref{gr:forest} with $q=\mu_nf$.  Taking absolute
values, grouping by $M$, and using
\eqref{gr:source-interpolation} proves
\[
 \|G_nf\|_\infty\leq A_n(M_n)m_f^\rho.
\]
The variable in \eqref{gr:An} is independent of the choice of $f$,
so this bound holds simultaneously for all admissible functions.
The tower property and \eqref{gr:Y-mean} give
$\sup_n\E A_n<\infty$, and Markov's inequality gives tightness.
\end{proof}

The estimate depends only on the finite map law and its spanning
trees.  In particular, it controls the full potential of a small
source, including any behavior on vertex sets that uniform
approximation of continuous functions might fail to detect.
